
\documentclass{article}
\usepackage{colortbl}
\usepackage{makecell}
\usepackage{multirow}
\usepackage{supertabular}

\begin{document}

\newcounter{utterance}

\centering \large Interaction Transcript for game `wordle\_withclue', experiment `high\_frequency\_words\_clue\_no\_critic', episode 7 with reka{-}flash{-}3.1{-}t0.0.
\vspace{24pt}

{ \footnotesize  \setcounter{utterance}{1}
\setlength{\tabcolsep}{0pt}
\begin{supertabular}{c@{$\;$}|p{.15\linewidth}@{}p{.15\linewidth}p{.15\linewidth}p{.15\linewidth}p{.15\linewidth}p{.15\linewidth}}
    \# & \multicolumn{2}{c}{Player} && \multicolumn{2}{c}{Game Master} \\
    \hline

    \theutterance \stepcounter{utterance}  
    & & & \multicolumn{4}{p{0.6\linewidth}}{
        \cellcolor[rgb]{0.9,0.9,0.9}{
            \makecell[{{p{\linewidth}}}]{
                \texttt{\tiny{[P1$\langle$GM]}}
                \texttt{You are a language wizard who likes to guess words by using the given rules.} \\
\\ 
\texttt{Welcome to Wordle! You have six attempts to guess the target word, a valid English word of five lowercase letters (a{-}z). Please use the tags "explanation:" and "guess:" to provide a concise explanation for each guess.} \\
\\ 
\texttt{To help you make an informed guess, you will receive a clue for the word, such as} \\
\texttt{clue: snowy white.} \\
\\ 
\texttt{Here is an example guess based on the clue:} \\
\texttt{explanation: In the fairy tail Snow White, the girl is killed because she eats a poisoned apple. And the word apple has 5 letters.} \\
\texttt{guess: apple} \\
\\ 
\texttt{After each guess, your answer will be validated, and you will receive feedback indicating which letters are correct (green), which letters are correct but in the wrong position (yellow), and which letters are incorrect (red). This feedback can be useful in determining which letters to include or exclude in your next guess.} \\
\\ 
\texttt{For example, the feedback for "apple" might be:} \\
\texttt{guess\_feedback: a<yellow> p<yellow> p<green> l<yellow> e<red>} \\
\\ 
\texttt{The explanation should contain details about how both the guess\_feedback and clue work together in tandem to formulate a new guess.} \\
\\ 
\texttt{Do not generate any other text.} \\
\\ 
\texttt{Let's begin with your first guess.} \\
\\ 
\texttt{clue: others'} \\
            }
        }
    }
     \\ \\

    \theutterance \stepcounter{utterance}  
    & \multicolumn{4}{p{0.6\linewidth}}{
        \cellcolor[rgb]{0.561,0.737,0.561}{
            \makecell[{{p{\linewidth}}}]{
                \texttt{\tiny{[P1$\rangle$GM]}}
                \texttt{<reasoning>} \\
\texttt{Let's tackle this Wordle with the clue "others'". The clue is possessive, so it's "others'". The target word is a five{-}letter English word. Let me think.} \\
\\ 
\texttt{First, "others'" is the possessive form of "others". So the word we're looking for might relate to the concept of "others" in some way. Maybe a word that means "another" or "additional"? Or perhaps a word that incorporates "others" as a part of it. Wait, but the target is five letters. Let's break it down.} \\
\\ 
\texttt{"Others" is six letters. The possessive is "others'". So maybe the target word is a part of "others'". Let's see: "others'" could be split into "other" (5 letters) plus an apostrophe and's', but the target is five letters. Wait, "other" is five letters. Could the target be "other"? But "other" is a valid word. Let me check if that's possible. But maybe the clue is more about the possessive form. Alternatively, maybe the word is "others" but that's six letters. Hmm.} \\
\\ 
\texttt{Alternatively, maybe the clue is pointing to a word that is related to the concept of "others" in a different way. For example, "another" is six letters. "Alone" is five letters, but that's the opposite. Wait, maybe "other" is the key here. Let's think: "other" is five letters. If the clue is "others'", maybe the target is "other" because "others'" is the possessive of "others", but "other" is a synonym. Alternatively, maybe the word is "yours" but that's five letters. Wait, "yours" is possessive but not directly related to "others'". Hmm.} \\
\\ 
\texttt{Wait, the clue is "others'", so maybe the target word is "other" (since "others" is plural, but "other" is singular). Let's check if "other" is a possible answer. But maybe there's another angle. Let's think of words that include "others" in some form. For example, "others" could be part of a longer word, but the target is five letters. Alternatively, maybe the word is "their" (5 letters), but that's possessive but not directly "others'". Alternatively, "whose"} \\
            }
        }
    }
    & & \\ \\

    \theutterance \stepcounter{utterance}  
    & & & \multicolumn{2}{p{0.3\linewidth}}{
        \cellcolor[rgb]{0.9,0.9,0.9}{
            \makecell[{{p{\linewidth}}}]{
                \texttt{\tiny{[GM$|$GM]}}
                \texttt{Error: The response should always start with the keyword 'explanation:'} \\
            }
        }
    }
    & & \\ \\

    \theutterance \stepcounter{utterance}  
    & & & \multicolumn{2}{p{0.3\linewidth}}{
        \cellcolor[rgb]{0.9,0.9,0.9}{
            \makecell[{{p{\linewidth}}}]{
                \texttt{\tiny{[GM$|$GM]}}
                \texttt{game\_result = ABORT} \\
            }
        }
    }
    & & \\ \\

\end{supertabular}
}

\end{document}
